\section*{Chapitre \ref{ch:nw:servers} --- Serveurs}

\begin{solution}{exo:nw:servers:1}
$2\,000 \times (0.8 \times 3/5 + 0.2) = \qty{1\,360}{\nano\second}$~; la limite est la partie liée à la mémoire, $0.2 \times 2\,000 =
\qty{400}{\nano\second}$.
\end{solution}

\begin{solution}{exo:nw:servers:2}
$15\,000/570 = 26.3$~: 26 serveurs (la place permet 42). $15\,000/1\,250 = 12$ serveurs de 2 unités (la place permet 21)~: l'alimentation est la contrainte dans les deux cas.
\end{solution}

\begin{solution}{exo:nw:servers:3}
Le boost, si le chemin critique tourne sur un ou deux cœurs et que le reste de la puce reste calme (sa marge thermique et électrique est ce qui permet à un
seul cœur de booster)~; la fréquence de base si de nombreux cœurs sont occupés en même temps, car le processeur ne peut alors pas maintenir le boost sur tous.
\end{solution}

\begin{solution}{exo:nw:servers:4}
Les SMI sont traitées par le firmware, en dehors du système d'exploitation, et ne peuvent pas être masquées~; aucun réglage du système d'exploitation ne peut donc les éloigner d'un cœur. La société
peut demander au fournisseur quelles sources de SMI peuvent être désactivées sans danger, choisir un firmware en comportant moins, les compter à l'aide des compteurs du processeur, et
mettre en corrélation les pauses inexpliquées avec elles.
\end{solution}

\begin{solution}{exo:nw:servers:5}
\qty{367}{\nano\second} si tout est lié au cœur, \qty{220}{\nano\second} à 60~\%, \qty{110}{\nano\second} à 30~\%.
\end{solution}

\begin{solution}{exo:nw:servers:6}
Un chemin critique qui utilise plusieurs cœurs à la fois (un feed handler, une stratégie, une passerelle sur des cœurs distincts) tourne à la fréquence de base d'un processeur standard
lorsqu'ils sont tous occupés, et une fréquence qui dépend de la charge fait varier sa latence avec la charge. Un \omterm{def:nw:servers:oc}{surcadençage} de tous les cœurs fait tourner chaque cœur à une fréquence fixe
de sorte que la latence du chemin ne dépend pas de ce qui tourne par ailleurs.
\end{solution}

\begin{solution}{exo:nw:servers:7}
En dessous de 25~\% de temps lié au cœur~: $1 - (0.25 \times 4.8/5.0 + 0.75) = 1~\%$. Avec \qty{20}{\kilo\watt}~: 35 serveurs 9175F (560 cœurs), 38 serveurs 6544Y
(608), 16 serveurs 9755 à deux sockets (4\,096 cœurs), 10 serveurs surcadencés (240).
\end{solution}

\begin{solution}{exo:nw:servers:8}
Si le chemin est surtout lié à la mémoire ou aux périphériques, une horloge plus rapide change peu de chose (\cref{prop:nw:servers:clock})~; et le nouveau serveur peut ne pas
avoir boosté (d'autres cœurs occupés, un profil d'alimentation, des états de veille profonds), ou le chemin tourne peut-être sur le mauvais socket. Il faut mesurer la part liée au cœur et
la fréquence réelle à laquelle le cœur a tourné avant d'incriminer le benchmark.
\end{solution}

\begin{solution}{pb:nw:servers:1}
\begin{enumerate}
\item \qty{1\,491}{\nano\second} (9175F à 5.0~GHz), \qty{1\,685}{\nano\second} (9755 à 4.1), \qty{1\,528}{\nano\second} (surcadencé, 4.8).
\item \qty{1\,659}{\nano\second}, \qty{2\,241}{\nano\second} et \qty{1\,528}{\nano\second}.
\item $2\,241 / 1\,491 = 1.50$.
\item $1/(1 - 0.6) = 2.5$~: il ne resterait au plus que les \qty{611}{\nano\second} liés à la mémoire.
\item 17, 8 et 5.
\item 272, 2\,048 et 120.
\item L'alimentation pour les trois (la place permettrait 42, 21 et 21).
\item Trois serveurs 9175F (48 cœurs chauds) et quatre serveurs 9755 (1\,024 cœurs)~: $3 \times 570 + 4 \times 1\,250 = 6\,710$~W, dans une seule baie.
\item Le profil d'alimentation et les états de veille (une machine qui se met en veille n'est pas celle qui a été testée), le turbo et le verrouillage de fréquence (la raison pour laquelle elle a été
  achetée), la vitesse de la mémoire et le signalement des erreurs ECC (mémoire surcadencée), et les sources de SMI.
\item Les pannes qui n'apparaissent que sous charge soutenue et sous chaleur~: stabilité marginale du \omterm{def:nw:servers:oc}{surcadençage}, erreurs mémoire, bridage thermique.
\item Il faut contrôler les ventilateurs et la boucle de refroidissement liquide, comparer avec l'enveloppe du \omterm{def:nw:servers:burnin}{test de déverminage} (burn-in), déplacer ses stratégies vers une machine de réserve avant qu'elle ne soit bridée ou ne tombe en panne,
  et ouvrir un dossier auprès du fournisseur.
\item Un thread frère partage le pipeline et les caches du cœur avec le thread chaud, ce qui ajoute de la contention et de la gigue (jitter)~; le laisser inactif réserve le cœur
  au thread chaud.
\item \textbf{Résultat nommé.} À 60~\% de temps lié au cœur, le chemin critique prend 1\,491, 1\,685 et \qty{1\,528}{\nano\second} sur le 9175F, le 9755 et le
  serveur surcadencé~; une baie de \qty{10}{\kilo\watt} en accueille 17, 8 et 5.
\item En dessous d'environ 6~\% de temps lié au cœur, même la fréquence la plus basse (2.7~GHz) reste à moins de 5~\% de la plus rapide.
\item En recherche, en simulation et dans les autres travaux de la société (le livre~15), dans un centre de données où l'énergie est bon marché, et non dans la cage de colocation.
\item Le chemin sur deux fréquences (ou sur deux cœurs de fréquences différentes), ajusté à la \cref{prop:nw:servers:clock}, et la fréquence réelle à laquelle le cœur a tourné.
\item Une fréquence stable garde la distribution de latence serrée~; une fréquence plus élevée mais variable ajoute de la variance aux moments où la machine est
  la plus sollicitée.
\item Les pannes relèvent du fournisseur, selon ses conditions~: les pièces de rechange, le support sur site et les délais de remplacement figurent dans le contrat (chapitre~15).
\item «~Les serveurs à 16 cœurs à 5~GHz donnent au chemin critique sa meilleure fréquence à la plus faible puissance~; une machine surcadencée ne vaut sa prime que si notre
  chemin critique utilise beaucoup de cœurs à la fois. Les serveurs à 128 cœurs ont leur place en recherche, pas dans la cage.~»
\item La part du temps du chemin qui est liée au cœur.
\end{enumerate}
\end{solution}

\begin{solution}{iq:nw:servers:1}
Un chemin critique tourne sur quelques cœurs~; ce qui compte est la vitesse de chacun, et une puce comportant moins de cœurs peut booster plus haut et garde plus de cache
et de puissance par cœur. Les nombreux cœurs ajoutent un débit que le chemin critique n'utilise pas, et des fréquences plus basses.
\iqlookfor{vitesse par cœur contre débit.}
\end{solution}

\begin{solution}{iq:nw:servers:2}
Profil d'alimentation à performance maximale, états C profonds désactivés, politique de turbo fixe, SMT désactivé ou inutilisé sur les cœurs chauds, mémoire à la vitesse nominale sans
entrelacement entre sockets, signalement des erreurs et sources de SMI réduits autant que le fournisseur le permet~; puis il faut les auditer avec les réglages du système d'exploitation.
\iqlookfor{chaque réglage avec le mécanisme de latence qu'il traite.}
\end{solution}

\begin{solution}{iq:nw:servers:3}
Une interruption traitée par le firmware, invisible du système d'exploitation et non masquable par lui, qui arrête le processeur pendant sa durée (de quelques microsecondes
à quelques millisecondes)~; l'isolation de cœurs ne peut pas la supprimer, de sorte qu'elle apparaît sous forme de pauses inexpliquées.
\iqlookfor{firmware, non masquable, non isolable.}
\end{solution}

\begin{solution}{iq:nw:servers:4}
Si le chemin critique est lié au cœur et utilise plusieurs cœurs à la fois. Il faut demander la fréquence sous notre charge soutenue, les tests de stabilité et le \omterm{def:nw:servers:burnin}{test de déverminage},
les modes de défaillance du refroidissement, la consommation électrique, la garantie et les pièces de rechange, les réglages du firmware et les sources de SMI, ainsi que des références d'utilisateurs comparables.
\iqlookfor{comportement soutenu et support, non le chiffre de pointe.}
\end{solution}

\begin{solution}{iq:nw:servers:5}
On mesure la part liée au cœur du chemin en l'exécutant à deux fréquences et en ajustant $t(f) = t_{\mathrm{ref}}(\phi f_{\mathrm{ref}}/f + 1 - \phi)$, puis
on prédit~; on confirme sur un serveur de prêt configuré comme en production.
\iqlookfor{un modèle issu de la mesure, puis un test.}
\end{solution}

\begin{solution}{iq:nw:servers:6}
L'alimentation d'abord~: on budgète chaque serveur à sa consommation réelle et on remplit la baie de serveurs du chemin critique, de commutateurs, de matériel de synchronisation et de capture~; tout ce qui
n'a pas besoin d'être proche de la plateforme de négociation va ailleurs. On garde une machine de réserve et de la marge pour la croissance.
\iqlookfor{l'alimentation comme contrainte déterminante et une règle claire sur ce qui appartient à la cage.}
\end{solution}
